\documentclass[12pt,a4paper]{article}
\usepackage{amsmath}
\usepackage{amssymb}
\usepackage{amsthm}
\usepackage{geometry}
\usepackage{algorithm}
\usepackage{algpseudocode}
\usepackage{graphicx}
\usepackage{hyperref}

\geometry{margin=1in}

\title{Mathematical Theory of Random Matrix Image Denoising}
\author{Matrix Analysis Lab}
\date{\today}

\newtheorem{theorem}{Theorem}
\newtheorem{definition}{Definition}

\begin{document}

\maketitle

\begin{abstract}
This document provides a comprehensive mathematical explanation of image denoising using Random Matrix Theory (RMT), specifically the Marchenko-Pastur (M-P) law and the Generalized Covariance method. We detail the noise model, eigenvalue-based filtering, and the role of SymPy in parameter optimization.
\end{abstract}

\tableofcontents
\newpage

\section{Introduction}

Image denoising is a fundamental problem in image processing. We treat a collection of $n$ grayscale images of size $H \times W$ as a random matrix and apply Random Matrix Theory to separate signal from noise.

\subsection{Problem Setup}

Given:
\begin{itemize}
    \item $n$ clean images: $\mathbf{X}_{\text{clean}} \in \mathbb{R}^{n \times H \times W}$
    \item Each image has $p = H \times W$ pixels
    \item Aspect ratio: $y = \frac{p}{n}$
\end{itemize}

Goal: Remove noise while preserving signal structure.

\section{Noise Model: Laplacian Distribution}

\subsection{Why Laplacian Noise?}

Unlike Gaussian noise, Laplacian noise has heavier tails, making it more realistic for modeling outliers in real-world images.

\subsection{Mathematical Definition}

The Laplacian (double exponential) distribution with scale parameter $b$ has probability density function:

\begin{equation}
p(x) = \frac{1}{2b} \exp\left(-\frac{|x|}{b}\right)
\end{equation}

Properties:
\begin{itemize}
    \item Mean: $\mu = 0$ (zero-mean noise)
    \item Variance: $\sigma^2 = 2b^2$
    \item Standard deviation: $\sigma = b\sqrt{2}$
\end{itemize}

\subsection{Adding Noise to Images}

For each pixel in each image, we add independent Laplacian noise:

\begin{equation}
\mathbf{X}_{\text{noisy}}[i,h,w] = \mathbf{X}_{\text{clean}}[i,h,w] + \epsilon_{ihw}
\end{equation}

where $\epsilon_{ihw} \sim \text{Laplace}(0, b)$ are i.i.d. random variables.

Default scale: $b = 0.1$

After adding noise, we clip to valid pixel range:
\begin{equation}
\mathbf{X}_{\text{noisy}} = \text{clip}(\mathbf{X}_{\text{clean}} + \mathbf{E}, 0, 1)
\end{equation}

where $\mathbf{E} \in \mathbb{R}^{n \times H \times W}$ contains i.i.d. Laplacian noise.

\section{Data Matrix Construction}

\subsection{Reshaping Images to Matrix Form}

We convert the batch of noisy images into a data matrix:

\begin{equation}
\mathbf{X} \in \mathbb{R}^{p \times n}
\end{equation}

where:
\begin{itemize}
    \item Each column $\mathbf{x}_i \in \mathbb{R}^p$ is one image (flattened)
    \item $p = H \times W$ (number of features/pixels)
    \item $n$ = number of images (samples)
\end{itemize}

Mathematically:
\begin{equation}
\mathbf{X}[:, i] = \text{flatten}(\mathbf{X}_{\text{noisy}}[i, :, :])
\end{equation}

\subsection{Centering the Data}

Compute column-wise mean and center:

\begin{equation}
\boldsymbol{\mu} = \frac{1}{n}\sum_{i=1}^n \mathbf{x}_i \in \mathbb{R}^p
\end{equation}

\begin{equation}
\mathbf{X}_c = \mathbf{X} - \boldsymbol{\mu} \mathbf{1}_n^T
\end{equation}

where $\mathbf{1}_n = (1, 1, \ldots, 1)^T \in \mathbb{R}^n$.

\section{Sample Covariance Matrix}

\subsection{Definition}

The sample covariance matrix is:

\begin{equation}
\mathbf{S} = \frac{1}{n}\mathbf{X}_c \mathbf{X}_c^T \in \mathbb{R}^{p \times p}
\end{equation}

This is a symmetric positive semi-definite matrix.

\subsection{Eigenvalue Decomposition}

Compute the eigendecomposition:

\begin{equation}
\mathbf{S} = \mathbf{V}\mathbf{\Lambda}\mathbf{V}^T
\end{equation}

where:
\begin{itemize}
    \item $\mathbf{V} = [\mathbf{v}_1, \mathbf{v}_2, \ldots, \mathbf{v}_p] \in \mathbb{R}^{p \times p}$ are eigenvectors
    \item $\mathbf{\Lambda} = \text{diag}(\lambda_1, \lambda_2, \ldots, \lambda_p)$ are eigenvalues
    \item Sorted in descending order: $\lambda_1 \geq \lambda_2 \geq \cdots \geq \lambda_p \geq 0$
\end{itemize}

\section{Marchenko-Pastur Law}

\subsection{Theoretical Foundation}

For a random matrix with i.i.d. entries of variance $\sigma^2$, as $p, n \to \infty$ with fixed ratio $y = p/n$, the empirical eigenvalue distribution of $\mathbf{S}$ converges to the Marchenko-Pastur distribution.

\begin{theorem}[Marchenko-Pastur Law]
The support of the limiting eigenvalue distribution is:
\begin{equation}
[\lambda_-, \lambda_+] = \left[\sigma^2(1-\sqrt{y})^2, \sigma^2(1+\sqrt{y})^2\right]
\end{equation}
\end{theorem}

\subsection{Key Insight}

\begin{itemize}
    \item Eigenvalues in $[\lambda_-, \lambda_+]$ correspond to \textbf{noise}
    \item Eigenvalues $> \lambda_+$ correspond to \textbf{signal}
\end{itemize}

\subsection{Estimating Noise Variance}

We estimate $\sigma^2$ from the bulk eigenvalues:

\begin{enumerate}
    \item Filter positive eigenvalues: $\{\lambda_i : \lambda_i > 10^{-10}\}$
    \item Take bottom 50\%: $\lambda_{\text{bulk}}$
    \item Compute minimum: $\lambda_{\min}^{\text{bulk}} = \min(\lambda_{\text{bulk}})$
    \item Estimate noise variance:
    \begin{equation}
    \hat{\sigma}^2 = \frac{\lambda_{\min}^{\text{bulk}}}{(1-\sqrt{y})^2}
    \end{equation}
\end{enumerate}

\subsection{Computing Thresholds}

Upper threshold (signal cutoff):
\begin{equation}
\lambda_+ = \hat{\sigma}^2(1 + \sqrt{y})^2
\end{equation}

Lower threshold (noise floor):
\begin{equation}
\lambda_- = \hat{\sigma}^2(1 - \sqrt{y})^2
\end{equation}

\subsection{Signal Subspace Identification}

Identify signal eigenvalues:
\begin{equation}
\mathcal{I}_{\text{signal}} = \{i : \lambda_i > \lambda_+ + \epsilon\}
\end{equation}

where $\epsilon = 10^{-9}$ is numerical tolerance.

Number of signal components:
\begin{equation}
m = |\mathcal{I}_{\text{signal}}|
\end{equation}

\subsection{Denoising via Projection}

Form signal subspace matrix:
\begin{equation}
\mathbf{\Gamma} = [\mathbf{v}_i : i \in \mathcal{I}_{\text{signal}}] \in \mathbb{R}^{p \times m}
\end{equation}

Project centered data onto signal subspace:
\begin{equation}
\tilde{\mathbf{X}}_c = \mathbf{\Gamma}\mathbf{\Gamma}^T \mathbf{X}_c
\end{equation}

Add back mean:
\begin{equation}
\tilde{\mathbf{X}} = \tilde{\mathbf{X}}_c + \boldsymbol{\mu}\mathbf{1}_n^T
\end{equation}

\subsection{Reconstruction}

Reshape each column back to image:
\begin{equation}
\tilde{\mathbf{X}}_{\text{clean}}[i,:,:] = \text{reshape}(\tilde{\mathbf{X}}[:,i], H, W)
\end{equation}

Clip to valid range:
\begin{equation}
\tilde{\mathbf{X}}_{\text{clean}} = \text{clip}(\tilde{\mathbf{X}}_{\text{clean}}, 0, 1)
\end{equation}

\section{Generalized Covariance Method}

\subsection{Motivation}

The standard M-P law assumes simple noise structure. The generalized method extends this to more complex covariance structures with parameters $\sigma^2$, $a$, and $\beta$.

\subsection{Generalized Support Formula}

The noise eigenvalue support is:
\begin{equation}
[\lambda_+^{\text{gen}}, \lambda_-^{\text{gen}}] = \left[\sigma^2 \max_{t>0} g(t), \sigma^2 \min_{t \in (-1/a, 0)} g(t)\right]
\end{equation}

where $g(t)$ is defined as:
\begin{equation}
g(t) = \frac{y\beta(a-1)t + (at+1)((y-1)t-1)}{(at+1)(t^2+t)}
\end{equation}

\subsection{Parameters}

\begin{itemize}
    \item $\sigma^2 > 0$: Noise variance scale
    \item $a > 1$: Spike parameter controlling eigenvalue spread
    \item $\beta \in (0, 1)$: Mixing parameter controlling distribution shape
    \item $y = p/n$: Aspect ratio (fixed by data)
\end{itemize}

\subsection{Computing Support Bounds}

\textbf{Step 1:} Find maximum of $g(t)$ for $t > 0$:
\begin{equation}
g_+ = \max_{t > 0} g(t) = -\min_{t > 0} [-g(t)]
\end{equation}

This is computed numerically via optimization.

\textbf{Step 2:} Find minimum of $g(t)$ for $t \in (-1/a, 0)$:
\begin{equation}
g_- = \min_{t \in (-1/a, 0)} g(t)
\end{equation}

\textbf{Step 3:} Multiply by noise variance:
\begin{equation}
\lambda_+^{\text{gen}} = \sigma^2 \cdot g_+, \quad \lambda_-^{\text{gen}} = \sigma^2 \cdot g_-
\end{equation}

Note: $\lambda_+^{\text{gen}} < \lambda_-^{\text{gen}}$ (inverted compared to M-P law).

\subsection{Signal Identification}

Signal eigenvalues are those exceeding the upper bound:
\begin{equation}
\mathcal{I}_{\text{signal}}^{\text{gen}} = \{i : \lambda_i > \lambda_-^{\text{gen}} + \epsilon\}
\end{equation}

The rest of the denoising process (projection, reconstruction) is identical to M-P law.

\section{Parameter Optimization with SymPy}

\subsection{The Optimization Problem}

We need to find optimal $(\sigma^2, a, \beta)$ such that the lower bound matches the observed minimum eigenvalue.

\textbf{Objective:} Minimize
\begin{equation}
\mathcal{L}(\sigma^2, a, \beta) = \left(\sigma^2 \cdot g_+ - \lambda_{\min}^{\text{bulk}}\right)^2
\end{equation}

subject to:
\begin{align}
\sigma^2 &> 0 \\
a &> 1 \\
\beta &\in (0, 1)
\end{align}

\subsection{How SymPy Works}

\textbf{SymPy} is a Python library for symbolic mathematics. Here's how it's used in optimization:

\subsubsection{Step 1: Symbolic Function Definition}

We define $g(t)$ symbolically (though in practice we use numerical evaluation):

\begin{verbatim}
from scipy.optimize import minimize_scalar, minimize

def g(t):
    numerator = y * beta * (a - 1) * t + (a*t + 1) * ((y-1)*t - 1)
    denominator = (a*t + 1) * (t**2 + t)
    return numerator / denominator
\end{verbatim}

\subsubsection{Step 2: Nested Optimization}

For each candidate $(\sigma^2, a, \beta)$:

\begin{enumerate}
    \item \textbf{Inner optimization:} Find $g_+$
    \begin{equation}
    g_+ = \max_{t \in [10^{-6}, 100]} g(t)
    \end{equation}

    Implemented as:
    \begin{verbatim}
result_max = minimize_scalar(
    lambda t: -g(t),
    bounds=(1e-6, 100),
    method='bounded'
)
g_plus = -result_max.fun
    \end{verbatim}

    \item \textbf{Compute lower bound:}
    \begin{equation}
    \text{lower\_bound} = \sigma^2 \cdot g_+
    \end{equation}

    \item \textbf{Evaluate objective:}
    \begin{equation}
    \mathcal{L} = (\text{lower\_bound} - \lambda_{\min}^{\text{bulk}})^2
    \end{equation}
\end{enumerate}

\subsubsection{Step 3: Outer Optimization}

Use L-BFGS-B algorithm to minimize $\mathcal{L}$:

\begin{verbatim}
# Initial guess
x0 = [1.0, 2.0, 0.5]  # [sigma2, a, beta]

# Bounds
bounds = [(0.01, 10.0), (1.01, 10.0), (0.01, 0.99)]

# Optimize
result = minimize(
    objective,
    x0,
    method='L-BFGS-B',
    bounds=bounds,
    options={'maxiter': 200, 'ftol': 1e-8}
)

sigma2_opt, a_opt, beta_opt = result.x
\end{verbatim}

\subsection{Why SymPy is Useful}

Although we use \texttt{scipy.optimize} for numerical optimization, SymPy could be used for:

\begin{enumerate}
    \item \textbf{Symbolic derivatives:} Compute $\frac{\partial g}{\partial t}$ analytically
    \begin{verbatim}
import sympy as sp
t = sp.Symbol('t')
g_sym = (y*beta*(a-1)*t + (a*t+1)*((y-1)*t-1)) / ((a*t+1)*(t**2+t))
g_prime = sp.diff(g_sym, t)
    \end{verbatim}

    \item \textbf{Critical point analysis:} Solve $\frac{\partial g}{\partial t} = 0$ exactly
    \begin{verbatim}
critical_points = sp.solve(g_prime, t)
    \end{verbatim}

    \item \textbf{Gradient computation:} For faster optimization
    \begin{equation}
    \nabla_{\theta} \mathcal{L} = 2(\sigma^2 g_+ - \lambda_{\min}^{\text{bulk}}) \cdot \nabla_{\theta}[\sigma^2 g_+]
    \end{equation}
    where $\theta = (\sigma^2, a, \beta)$.
\end{enumerate}

\subsection{Optimization Algorithm Summary}

\begin{algorithm}
\caption{Generalized Parameter Optimization}
\begin{algorithmic}[1]
\State \textbf{Input:} Eigenvalues $\{\lambda_i\}$, aspect ratio $y$
\State Filter positive eigenvalues: $\lambda_{\text{pos}}$
\State Extract bottom 50\%: $\lambda_{\text{bulk}}$
\State $\lambda_{\min}^{\text{bulk}} \gets \min(\lambda_{\text{bulk}})$
\State Initialize: $\sigma^2_0 = 1.0$, $a_0 = 2.0$, $\beta_0 = 0.5$
\For{iteration $= 1, 2, \ldots, 200$}
    \State Evaluate $g_+(\sigma^2, a, \beta)$ via inner optimization
    \State Compute $\mathcal{L} = (\sigma^2 \cdot g_+ - \lambda_{\min}^{\text{bulk}})^2$
    \State Update $(\sigma^2, a, \beta)$ using L-BFGS-B step
    \If{$|\mathcal{L}| < 10^{-8}$}
        \State \textbf{break} (converged)
    \EndIf
\EndFor
\State \textbf{Return:} $(\sigma^2_{\text{opt}}, a_{\text{opt}}, \beta_{\text{opt}})$
\end{algorithmic}
\end{algorithm}

\section{Complete Denoising Pipeline}

\subsection{Algorithm Overview}

\begin{algorithm}
\caption{Random Matrix Image Denoising}
\begin{algorithmic}[1]
\State \textbf{Input:} Clean images $\mathbf{X}_{\text{clean}} \in \mathbb{R}^{n \times H \times W}$, noise scale $b$, method
\State \textbf{Output:} Denoised images $\tilde{\mathbf{X}}_{\text{clean}}$
\State
\State // \textbf{Step 1: Add Laplacian Noise}
\State $\mathbf{E} \sim \text{Laplace}(0, b)^{n \times H \times W}$
\State $\mathbf{X}_{\text{noisy}} \gets \text{clip}(\mathbf{X}_{\text{clean}} + \mathbf{E}, 0, 1)$
\State
\State // \textbf{Step 2: Reshape to Matrix}
\State $\mathbf{X} \gets \text{images\_to\_matrix}(\mathbf{X}_{\text{noisy}})$ \Comment{Shape: $(p, n)$}
\State $p \gets H \times W$, $y \gets p/n$
\State
\State // \textbf{Step 3: Center Data}
\State $\boldsymbol{\mu} \gets \frac{1}{n}\mathbf{X}\mathbf{1}_n$
\State $\mathbf{X}_c \gets \mathbf{X} - \boldsymbol{\mu}\mathbf{1}_n^T$
\State
\State // \textbf{Step 4: Compute Covariance}
\State $\mathbf{S} \gets \frac{1}{n}\mathbf{X}_c \mathbf{X}_c^T$
\State
\State // \textbf{Step 5: Eigendecomposition}
\State $(\mathbf{\Lambda}, \mathbf{V}) \gets \text{eigh}(\mathbf{S})$
\State Sort eigenvalues/vectors in descending order
\State
\State // \textbf{Step 6: Estimate Noise and Find Signal}
\If{method = "M-P Law"}
    \State $\hat{\sigma}^2 \gets \frac{\min(\text{bottom 50\%}(\lambda_{\text{pos}}))}{(1-\sqrt{y})^2}$
    \State $\lambda_+ \gets \hat{\sigma}^2(1+\sqrt{y})^2$
    \State $\mathcal{I}_{\text{signal}} \gets \{i : \lambda_i > \lambda_+ + 10^{-9}\}$
\ElsIf{method = "Generalized Covariance"}
    \State $(\sigma^2, a, \beta) \gets \text{optimize\_parameters}(\{\lambda_i\}, y)$
    \State $g_+ \gets \max_{t>0} g(t; a, \beta, y)$
    \State $g_- \gets \min_{t \in (-1/a,0)} g(t; a, \beta, y)$
    \State $\lambda_+^{\text{gen}} \gets \sigma^2 \cdot g_+$
    \State $\lambda_-^{\text{gen}} \gets \sigma^2 \cdot g_-$
    \State $\mathcal{I}_{\text{signal}} \gets \{i : \lambda_i > \lambda_-^{\text{gen}} + 10^{-9}\}$
\EndIf
\State
\State // \textbf{Step 7: Project onto Signal Subspace}
\State $\mathbf{\Gamma} \gets [\mathbf{v}_i : i \in \mathcal{I}_{\text{signal}}]$
\State $\tilde{\mathbf{X}}_c \gets \mathbf{\Gamma}\mathbf{\Gamma}^T \mathbf{X}_c$
\State $\tilde{\mathbf{X}} \gets \tilde{\mathbf{X}}_c + \boldsymbol{\mu}\mathbf{1}_n^T$
\State
\State // \textbf{Step 8: Reshape Back to Images}
\State $\tilde{\mathbf{X}}_{\text{clean}} \gets \text{matrix\_to\_images}(\tilde{\mathbf{X}}, H, W)$
\State $\tilde{\mathbf{X}}_{\text{clean}} \gets \text{clip}(\tilde{\mathbf{X}}_{\text{clean}}, 0, 1)$
\State
\State \textbf{return} $\tilde{\mathbf{X}}_{\text{clean}}$
\end{algorithmic}
\end{algorithm}

\section{Summary of Key Formulas}

\subsection{Noise Model}
\begin{equation}
\epsilon \sim \text{Laplace}(0, b), \quad \mathbf{X}_{\text{noisy}} = \text{clip}(\mathbf{X}_{\text{clean}} + \mathbf{E}, 0, 1)
\end{equation}

\subsection{Marchenko-Pastur Method}
\begin{align}
\hat{\sigma}^2 &= \frac{\lambda_{\min}^{\text{bulk}}}{(1-\sqrt{y})^2} \\
\lambda_+ &= \hat{\sigma}^2(1+\sqrt{y})^2 \\
\mathcal{I}_{\text{signal}} &= \{i : \lambda_i > \lambda_+\}
\end{align}

\subsection{Generalized Covariance Method}
\begin{align}
g(t) &= \frac{y\beta(a-1)t + (at+1)((y-1)t-1)}{(at+1)(t^2+t)} \\
g_+ &= \max_{t>0} g(t), \quad g_- = \min_{t \in (-1/a,0)} g(t) \\
\lambda_+^{\text{gen}} &= \sigma^2 g_+, \quad \lambda_-^{\text{gen}} = \sigma^2 g_- \\
\mathcal{I}_{\text{signal}} &= \{i : \lambda_i > \lambda_-^{\text{gen}}\}
\end{align}

\subsection{Denoising Projection}
\begin{align}
\mathbf{\Gamma} &= [\mathbf{v}_i : i \in \mathcal{I}_{\text{signal}}] \\
\tilde{\mathbf{X}} &= \mathbf{\Gamma}\mathbf{\Gamma}^T (\mathbf{X} - \boldsymbol{\mu}\mathbf{1}_n^T) + \boldsymbol{\mu}\mathbf{1}_n^T
\end{align}

\section{Computational Complexity}

\subsection{Time Complexity}

For $n$ images of size $H \times W$ with $p = HW$:

\begin{itemize}
    \item \textbf{Adding noise:} $O(nHW) = O(np)$
    \item \textbf{Reshaping:} $O(np)$
    \item \textbf{Centering:} $O(np)$
    \item \textbf{Covariance:} $O(np^2)$
    \item \textbf{Eigendecomposition:} $O(p^3)$ (dominant cost)
    \item \textbf{Projection:} $O(npm)$ where $m$ = number of signal components
    \item \textbf{Optimization (Generalized):} $O(K \cdot I)$ where $K$ = iterations, $I$ = inner optimization cost
\end{itemize}

\textbf{Total:} $O(p^3 + np^2)$ dominated by eigendecomposition.

\subsection{Space Complexity}

\begin{itemize}
    \item \textbf{Data matrix:} $O(np)$
    \item \textbf{Covariance matrix:} $O(p^2)$
    \item \textbf{Eigenvectors:} $O(p^2)$
\end{itemize}

\textbf{Total:} $O(p^2 + np)$

\section{Conclusion}

This document presented the complete mathematical theory behind random matrix image denoising:

\begin{enumerate}
    \item \textbf{Laplacian noise model} with heavy-tailed distribution
    \item \textbf{Marchenko-Pastur law} for simple i.i.d. noise
    \item \textbf{Generalized covariance method} for complex noise structures
    \item \textbf{SymPy-based optimization} for automatic parameter tuning
\end{enumerate}

The key insight is that Random Matrix Theory provides a principled way to separate signal from noise in the eigenvalue spectrum, enabling effective denoising through subspace projection.

\end{document}
